\newpage
\section{FPGA Backend}

The FPGA backend targets Xilinx Zynq devices using the PYNQ framework. It offloads the most computationally intensive operations—field updates, interference computation, and pattern detection—to custom hardware overlays, enabling true parallel execution. This backend demonstrates APEIRON-Core's ability to integrate radically different computational fabrics under a unified API.

\subsection{Hardware Overlay Architecture}

The FPGA is programmed with a custom bitstream that instantiates three primary IP blocks: \texttt{continuous\_fields}, \texttt{interference\_matrix}, and \texttt{stability\_detector}. The bitstream was synthesized for the Xilinx ZCU104 board (Zynq UltraScale+ XCZU7EV) using Vitis HLS 2023.2. The design operates at 200\,MHz and consumes approximately 38\% of the device's LUTs, 22\% of BRAM, and 15\% of DSP slices. The overlay is accessible via the PYNQ \texttt{Overlay} class, and all register interfaces are documented in the accompanying hardware user guide~\cite{beniers2025fpga}.

The overlay communicates with the ARM processor via AXI DMA channels. Interrupts are used to signal completion, though a polling fallback is provided for environments where interrupts are unreliable.

\subsection{Fixed-Point Conversion for AXI Registers}

A critical challenge in FPGA acceleration is the impedance mismatch between Python's floating-point numbers and the integer registers exposed by AXI interfaces. APEIRON-Core solves this by converting floating-point parameters (e.g., the time step \(\Delta t\)) to Q16.16 fixed-point format before writing to hardware registers:

\begin{lstlisting}[language=Python]
FIXED_POINT_SCALE = 65536  # 2^16
def _float_to_fixed(self, value: float) -> int:
    scale = self.config.get('fixed_point_scale', self.FIXED_POINT_SCALE)
    return int(round(value * scale))
\end{lstlisting}

The \texttt{round} operation minimizes quantization error. For parameters that must be represented exactly (e.g., floating-point constants), the backend also supports direct IEEE 754 bit-level transfer by reinterpreting the float as a 32-bit integer.

\subsection{DMA Buffer Management}

Field data is transferred between the ARM processor and the FPGA fabric via DMA. The backend uses PYNQ's \texttt{allocate} function to create physically contiguous buffers in the FPGA's address space. Each buffer is assigned a unique ID, and a hash-based cache avoids redundant allocation for identical fields. The \texttt{field\_update} method writes the IDs of the participating buffers and the fixed-point \(\Delta t\) to the hardware registers, then starts the IP block:
\newpage
\begin{lstlisting}[language=Python, caption={Initiating an FPGA field update.}]
self.field_engine.write(1, field_id)
self.field_engine.write(2, verleden_id)
self.field_engine.write(3, heden_id)
self.field_engine.write(4, toekomst_id)
self.field_engine.write(5, dt_fixed)
self.field_engine.start()
\end{lstlisting}

After completion, the updated field buffers are read back and the cached NumPy arrays are refreshed.

\subsection{CPU Fallback}

When PYNQ is not installed or no FPGA board is detected, the backend automatically falls back to a CPU implementation. This ensures that higher APEIRON layers can run unmodified in development environments, with hardware acceleration transparently engaged when deployed on a compatible device.